\chapter{Antibiotic Sensitivity Test}

\section*{What Is This Test?}
An antibiotic sensitivity test, also called antimicrobial susceptibility testing (AST), measures whether a cultured bacterium or fungus is inhibited by selected antimicrobial drugs. It helps the clinical team choose treatment that is likely to work against the organism recovered from the patient's specimen.

\section*{Alternative Names}
\begin{itemize}
  \item Antimicrobial Susceptibility Testing (AST)
  \item Antibiogram Testing
  \item Culture and Sensitivity (C\&S)
  \item Antibiotic Susceptibility Test
\end{itemize}

\section*{Principle}
The laboratory exposes a standardized isolate to antimicrobial agents. Disk diffusion reports the zone of inhibition, while broth microdilution and automated systems determine a minimum inhibitory concentration (MIC). The result is interpreted using organism-, drug-, and site-specific breakpoints as susceptible, intermediate or susceptible with increased exposure, or resistant.

\section*{Why This Test Is Ordered}
\begin{itemize}
  \item To select or narrow antimicrobial treatment after a clinically significant organism is cultured.
  \item To investigate treatment failure, recurrent infection, or suspected multidrug resistance.
  \item To support infection-control surveillance and local antibiograms.
\end{itemize}

\section*{Primary vs Secondary Test}
AST is a secondary test performed on a clinically significant isolate from a culture. It does not replace a good-quality specimen, organism identification, or clinical assessment. Molecular resistance-gene testing may provide earlier information but does not always predict susceptibility to every drug.

\section*{How This Test Is Performed}
The patient provides the specimen appropriate to the suspected infection, such as blood, urine, sputum, wound material, or cerebrospinal fluid. The laboratory cultures and identifies the organism, then tests selected antimicrobial agents. Results may take several hours after growth is detected or longer for slow-growing organisms.

\section*{What Preparation Is Needed}
Preparation depends on the specimen. When safe and feasible, cultures are collected before antimicrobial therapy begins. Do not delay urgent treatment to obtain a culture.

\section*{Contraindications and Precautions}
Contamination, colonization, mixed growth, poor specimen collection, and prior antimicrobial exposure can make results misleading. A susceptible result does not guarantee clinical success if drug penetration, dose, adherence, source control, or host factors are unfavorable.

\section*{Understanding Results}
Interpret the category using the laboratory's current breakpoint standard. MIC values should not be compared across different organisms or drugs without their corresponding breakpoints. Treatment should be narrowed when appropriate, while severe infection may require empiric coverage before AST is available.

\section*{Follow-up Tests}
Follow-up may include repeat culture for persistent or recurrent infection, blood cultures for suspected bloodstream infection, susceptibility testing of unusual isolates, and source-control evaluation or imaging when clinically indicated.

\section*{References}
\begin{enumerate}
  \item MedlinePlus. Antibiotic Sensitivity Test.
  \item Clinical and Laboratory Standards Institute. Performance Standards for Antimicrobial Susceptibility Testing.
  \item World Health Organization. Global Antimicrobial Resistance and Use Surveillance System (GLASS).
\end{enumerate}

\clearpage
\section*{Understanding Results and Follow-up}
The laboratory report should identify the specimen, method, units, reference interval or decision limit, and whether the result is positive, negative, abnormal, or inconclusive. Reference intervals vary with age, sex, pregnancy, specimen type, assay, and laboratory; a result outside a range is not automatically a diagnosis, and a result within a range does not always exclude disease. Discuss unexpected results with the ordering clinician, who can decide whether repeat or confirmatory testing, treatment, or referral is appropriate. Do not change medicines or delay urgent care based on an isolated result.


\section*{Regulatory Status and Authorization Date}
This chapter describes a test category, not a specific commercial product. FDA, CDSCO, EU IVDR, and MHRA approval or registration dates therefore cannot be assigned without the assay manufacturer, product name, instrument, and intended use. For laboratory-developed tests, an individual agency approval date may not exist. See the Preface for the interpretation of regulatory status.

\clearpage
